% 01-defi.tex — Le défi et la théorie du cours (vitesse minimale, tension maximale).

\section{Le défi et la théorie du cours}

Le mode jeu de l'exercice du looping propose une expérience que chacun connaît :
faire tourner un seau d'eau dans un plan vertical sans se mouiller. Une barre
tourne autour d'un pivot ; une corde pend à son bout et porte un seau rempli
d'eau. L'étudiant règle une seule grandeur, tirée au sort à chaque défi : la
vitesse \(v_0\) du seau au point bas, la longueur \(L\) de la corde ou la masse
\(m\) du seau rempli. Le défi est réussi si le seau fait deux tours sans perdre
une goutte ni casser la corde (\cref{fig-captures}).

Dans le cours, le seau est un point matériel au bout d'une corde sans masse,
à la distance \(R\) du pivot. La conservation de l'énergie donne la vitesse à
l'angle \(\theta\) compté depuis le point bas,
\(v^2 = v_0^2 - 2 g R (1 - \cos(\theta))\), et la deuxième loi de Newton selon le
rayon donne la tension \cite{formulaire_epl1201} :
\begin{equation}
  T = m \left(\frac{v^2}{R} + g \cos(\theta)\right)
    = m \left(\frac{v_0^2}{R} + g \left(3 \cos(\theta) - 2\right)\right).
  \label{eq-tension}
\end{equation}
L'eau reste au fond tant que la corde tire (\(T > 0\)), et c'est au sommet
(\(\theta = \pi\)) que c'est le plus difficile :
\begin{equation}
  T_s = m \left(\frac{v_0^2}{R} - 5 g\right) \geq 0
  \quad \Longleftrightarrow \quad v_0 \geq \sqrt{5 g R} .
  \label{eq-sommet}
\end{equation}
C'est la condition du wagon dans le looping, \(h \geq 5R/2\), avec
\(v_0^2 = 2 g h\). La corde casse si la tension dépasse sa tension de rupture
\(T_{\max}\) ; la plus grande tension est au point bas,
\(m (v_0^2/R + g)\). L'étudiant calcule ces deux bornes. Ce rapport explique
comment le jeu simule le seau, la corde et l'eau, puis compare ce que la
simulation obtient à ce que ces formules prévoient.
